\section{Appendix}
\label{app:appendix}

\subsection{Reproducibility: Experiment Runners}
\label{app:repro}
The reported experiments are executed from the Battery-TA-RWARE repository using the objective-family runners for centralized and shared-context planning.
Earlier exploratory scripts are not part of the reported comparison.
The examples below show the command structure used by the current protocol; the seed, prompt format, model subset, and result directory can be changed for a particular session.

\begin{lstlisting}[caption=Run Objective 1 with the shared-context planner., label=listing:run_obj1_shared]
python Battery-TA-RWARE/scripts/run_obj1_shared_context_llm.py \
  --prompt_format language \
  --max_steps 0 \
  --max_busy_steps_for_replan 10 \
  --max_inactivity_steps 500 \
  --disable_support_needed_soon \
  --find_path_agent_aware_always \
  --seed 0
\end{lstlisting}

\begin{lstlisting}[caption=Run Objective 1 with the centralized planner., label=listing:run_obj1_centralized]
python Battery-TA-RWARE/scripts/run_obj1_centralized_llm.py \
  --prompt_format language \
  --max_steps 0 \
  --max_busy_steps_for_replan 10 \
  --max_inactivity_steps 500 \
  --disable_support_needed_soon \
  --find_path_agent_aware_always \
  --seed 0
\end{lstlisting}

\begin{lstlisting}[caption=Run Objective 2 with the shared-context planner., label=listing:run_obj2_shared]
python Battery-TA-RWARE/scripts/run_obj2_shared_context_llm.py \
  --prompt_format language \
  --max_steps 0 \
  --max_busy_steps_for_replan 10 \
  --max_inactivity_steps 500 \
  --disable_support_needed_soon \
  --find_path_agent_aware_always \
  --seed 0
\end{lstlisting}

\begin{lstlisting}[caption=Run Objective 2 with the centralized planner., label=listing:run_obj2_centralized]
python Battery-TA-RWARE/scripts/run_obj2_centralized_llm.py \
  --prompt_format language \
  --max_steps 0 \
  --max_busy_steps_for_replan 10 \
  --max_inactivity_steps 500 \
  --disable_support_needed_soon \
  --find_path_agent_aware_always \
  --seed 0
\end{lstlisting}

\begin{lstlisting}[caption=Run Objective 3 with the shared-context planner., label=listing:run_obj3_shared]
python Battery-TA-RWARE/scripts/run_obj3_shared_context_llm.py \
  --prompt_format language \
  --max_steps 0 \
  --max_busy_steps_for_replan 10 \
  --max_inactivity_steps 500 \
  --disable_support_needed_soon \
  --find_path_agent_aware_always \
  --seed 0
\end{lstlisting}

\begin{lstlisting}[caption=Run Objective 3 with the centralized planner., label=listing:run_obj3_centralized]
python Battery-TA-RWARE/scripts/run_obj3_centralized_llm.py \
  --prompt_format language \
  --max_steps 0 \
  --max_busy_steps_for_replan 10 \
  --max_inactivity_steps 500 \
  --disable_support_needed_soon \
  --find_path_agent_aware_always \
  --seed 0
\end{lstlisting}

The same commands can be run with \nolinkurl{--prompt_format json} to execute the structured prompt variant.
Objective 1 also supports \nolinkurl{--selected_models}.
For Objectives 2 and 3, the corresponding single-model filter is \nolinkurl{--only_model}.
The seed is logged for traceability.

\subsection{Result Directories}
\label{app:result_dirs}
The objective-family runners write session outputs under the following default directories:
\begin{itemize}
    \item \nolinkurl{results/obj1_shared_context_llm_experiments}
    \item \nolinkurl{results/obj1_centralized_llm_experiments}
    \item \nolinkurl{results/obj2_shared_context_llm_experiments}
    \item \nolinkurl{results/obj2_centralized_llm_experiments}
    \item \nolinkurl{results/obj3_shared_context_llm_experiments}
    \item \nolinkurl{results/obj3_centralized_llm_experiments}
\end{itemize}
These raw step-level result directories are retained separately because their size is unsuitable for the source repository; the repository contains the experiment implementation and runner commands.
Each session stores aggregate summaries together with per-run JSON summaries and detailed logs.
The aggregate analysis uses the retained runs from these outputs and follows the protocol described in Section~\ref{sec:method_scenarios_models_metrics}.

\subsection{Prompt Excerpt}
\label{app:prompt_excerpt}
The full prompts are generated programmatically from the simulator state.
The excerpt below is illustrative and shows the action-id grounding and output contract used by the natural-language shared-context prompt.

\begin{lstlisting}[caption=Illustrative shared-context natural-language prompt excerpt., label=listing:prompt_excerpt]
Objective:
Maximize useful warehouse progress under the current objective.

Current agent:
agent_0, role=AGV, position=(8, 3), battery=72, carrying=None

Shared context:
- Requested shelf 24 is at position (11, 5), action id 50.
- Picker support is required at shelf action id 50.
- Charging station action id 7 is currently unoccupied.

Candidate actions:
- action=0, type=NOOP, distance=0, instruction=wait only if blocked
- action=50, type=SHELF, distance=8, instruction=go to requested shelf
- action=7, type=CHARGING, distance=11, instruction=use if battery is unsafe

Output contract:
Return only:
Action: <integer action id>
Steps: <integer hold duration>
Reason: <short reason>
\end{lstlisting}

For the centralized planner, the prompt contains the shared warehouse context plus one planning block for each agent being replanned.
The centralized response must contain one action for each listed agent.
For the JSON prompt variant, the same information is represented as a structured object and the model is constrained to return a machine-readable action object or action list.

\subsection{Run Summary Fields}
\label{app:run_summary_fields}
The aggregate CSV files use the following run-level fields:
\begin{lstlisting}[caption=Run-summary CSV fields used for aggregation., label=listing:summary_fields]
Objective
Arch
Prompt
Model
Scenario
Seed
Run
llm_calls
json_generation_failures
llm_missing_or_invalid_actions
charging_station_utilization_rate
elapsed_seconds
steps to first delivery
total deliveries
total steps
deliveries_per_1000_steps
llm_calls_per_1000_steps
llm_calls_per_delivery
invalid_action_rate
json_failure_rate
summary_json_path
\end{lstlisting}

The main analysis uses total deliveries, \gls{llm} calls, deliveries per 1000 \gls{llm} calls, and the descriptive battery-aware productivity score defined in Section~\ref{sec:method_scenarios_models_metrics}.
The remaining fields are retained for filtering, traceability, and diagnostic interpretation.
Parsing and invalid-action fields support diagnostic tracing of malformed outputs and configuration errors.
The \nolinkurl{summary_json_path} field links each aggregate row to the corresponding detailed run summary.
